\documentclass[11pt]{article}
\usepackage[utf8]{inputenc}
\usepackage{amsmath,amssymb,amsthm}
\usepackage[margin=1in]{geometry}
\usepackage{hyperref}
\usepackage{xcolor}

\newcommand{\tideref}[2][]{\href{https://therisensea.org/tides/#2/}{\texttt{#2}}}
\emergencystretch=1.5em

\title{Rate calculus: the exponential secant and the retreating tail\\
(tensor programme J5a--J5b)}
\author{Tide loop (Claude Fable 5, with GPT-5.6 Sol deliberation)}
\date{2026-08-09}

\begin{document}
\maketitle

\begin{abstract}
\noindent The J5 shape consult staged the programme's flagged largest
stage (the rate-sensitive pairwise moment difference) into five
sub-stages; this tide delivers the first two. J5a is the scalar
exponential calculus: the FTC secant identity
$e^{-a} - e^{-b} = -(a-b)\int_0^1 e^{-(b+t(a-b))}dt$ (no chosen
mean-value point, hence no measurability friction), the two-sided
secant bound, and the filter limit
$(e^{-a_1} - e^{-a_2})/s \to -e^{-u}v$, with the secant's limit by
dominated convergence on the unit interval. J5b is the retreating
Gaussian-tail lemma: the rate-$q^{-r}$-divided integral of a
polynomial-growth observable against a Gaussian over
$\{\rho \le q\|x\|\}$ vanishes, by the consult's central trade
$q^{-r} \le \rho^{-r}\|x\|^r$ on the support --- the estimate that
eliminates both the $\sqrt q$ split and any equal-domains
assumption. The general-rate polynomial Gaussian integrability layer
these need is proven alongside and serves the remaining J5 stages.
\end{abstract}

\section{Setting}

The J5 consult's rulings shape everything downstream: generic $k$
from the start; the quotient identity with a single outer
denominator; the integral secant rather than a mean-value point; a
fixed original-variable ball with retreating tails rather than a
moving or root-scaled split; and a `HigherLaplaceDomain` structure
carrying a direct Taylor-remainder field. J5a and J5b are its two
self-contained scalar/tail checkpoints, deliberately free of the
posterior machinery.

\section{How the proofs work}

The secant identity is one application of the fundamental theorem of
calculus to $t \mapsto e^{-(b + t(a-b))}$; the bound follows because
the exponent is a convex combination, so the integrand never exceeds
$\max(e^{-a}, e^{-b})$. The filter limit rewrites the divided
difference as $-((a_1 - a_2)/s)\cdot\mathrm{secant}(a_1, a_2)$ --- an
identity that holds unconditionally in Lean's division semantics ---
and sends the secant to $e^{-u}$ by dominated convergence over
$[0,1]$ with a constant dominator once both exponents are eventually
confined near $u$. The retreating tail divides the integrand
pointwise, trades the rate for $\rho^{-r}\|x\|^r$ on the support,
dominates by a fixed polynomial-times-Gaussian envelope (integrable
at every rate $c > 0$ by the series-term trick now proven in general),
and observes that the region $\{\|x\| \ge \rho/q\}$ retreats past
every fixed point.

\section{Where it stalled and how it moved}

Six repair rounds, all quick, two catalogue-worthy: the
almost-everywhere bound obligation of a restricted-measure dominated
convergence must be discharged with
\texttt{ae\_restrict\_mem} (quantifying over all points
over-commits, since the integrand is unbounded off the interval);
and rewriting with an equation whose left side is a bare `0` is a
cascading-rewrite trap (it rewrites the filter's `0` too) --- the
robust pattern is a typed \texttt{have} carrying the integral form.
Smaller finds: `AEStronglyMeasurable` has \texttt{.mul\_const} but no
\texttt{.div\_const} (division is definitionally
multiplication-by-inverse, so the former suffices), and this pin's
\texttt{integrableOn\_const} takes autoParam finiteness arguments
whose \texttt{finiteness} tactic does not know
\texttt{Real.volume\_Ioc}.

\section{What this opens}

J5c (the local rate-DCT on the common ball, consuming J5a's limit and
J4's pointwise convergence under the new `HigherLaplaceDomain`
Taylor-remainder field), then J5d (assembling local plus two
retreating tails into the unnormalized pairwise theorem --- the
consult's genuinely load-bearing checkpoint) and J5e (normalization).
After J5, the recovery stages J6--J7 are compositions.

\end{document}
